% !TEX encoding = UTF-8 Unicode
% !TEX root =  ../Main.tex

\chapter{IST-Analyse}\label{cha:ist}
Dieses Kapitel analysiert den aktuellen Prozess der Antragsstellung zur Erstattung des Semesterticketbeitrags durch eine studierende Person der \gls{hwr} beim \gls{asta} der \gls{hwr}. Da Ziel ist es durch diese Analyse Schwachstellen in den Prozess zu erkennen, um diese in einem späteren Kapitel mithilfe eines Softwareprototyps zu digitalisieren, optimieren und/oder automatisieren. Optimieren meint in diesem Fall die Antragsbearbeitung durch die Mitarbeitenden des Sozialbüros zu erleichtern und weniger zeitaufwändig zu gestalten. 

\section{Ausgangssituation}\label{sec:Ausgangssituation}
Der \gls{asta} der \gls{hwr} schließt in der Regel mit einem Verkehrsanbieter einen Vertrag über den Erwerb des Semestertickets durch Studierende ab. Er vertritt dabei die Studierendenschaft. Der aktuelle Vertrag wurde am 13.02.2025 mit der S-Bahn Berlin GmbH sowie der Verkehrsbund Berlin-Brandenburg GmbH geschlossen. Der Vertrag regelt in §1, dass alle ordentlich immatrikulierten Studierenden der \gls{hwr} verpflichtet sind das Deutschlandsemesterticket zu beziehen. Davon ausgeschlossen sind Gasthörer:innen und Zweithörer:innen, Studierende in Abend-, Online- und Fernstudiengängen ohne Präsenzpflicht, Studierende in berufsbegleitenden und weiterbildenden Studiengängen bei denen das Studium zeitlich unterwiegt, und Studierende die nicht der Studierendenschaft angehören. 
Des weiteren können Studierende aufgrund der folgenden Gründe eine Rückerstattung des Semesterticketbeitrags beantragen: 
\begin{itemize}
  \item Urlaubssemester
  \item Auslandssemester oder Praxissemester mit einer Aufenthaltsdauer von mehr als drei zusammenhängenden Monaten außerhalb des Geltungsbereichs
  \item Doppeltimmatrikulation an zwei Hochschulen mit Semesterticketpflicht
  \item Verspätete Immatrikulation nach nach mehr als einem Monat nach Semesterbeginn
  \item Exmatrikulation im laufenden Semester
  \item Chronische Erkrankung oder Behinderung
\end{itemize}

\section{Prozessbeschreibung}\label{sec:Prozessbeschreibung}
Um eine Rückserstattung zu beantragen müssen die Studierenden einen Antrag ausfüllen und diesen mitsamt ihrer Studienbescheinigung und weiter benötigter Nachweise per E-Mail an das Sozialbüro des \gls{asta} senden. Die Mitarbeitenden überprüfen die Dokumente dann auf ihre Korrektheit und Vollständigkeit. Wenn nicht alle Dokumente eingereicht wurden oder diese fehlerhaft sind, müssen diese erneut bei den Studierenden angefragt werden. Sind alle Dokumente korrekt und vollständig, werden die Daten in eine Excel Tabelle eingetragen und die Dokumente mit einer Prozessnummer versehen, um sie dem jeweiligen Antrag genau zuordnen zu können. Die Excel Tabelle sowie die Dokumente werden in einem privaten Google-Drive-Verzeichnis abgelegt. Nach dieser Erstprüfung wird die E-Mail der/des Studierenden mit einem Tag versehen und zur Zweitprüfung in das E-Mail-Postfach einer/eines anderen Mitarbeitenden verschoben. Der/die Zweitprüfende sucht dem Namen zugeordnet die Prozessnummer aus der Excel-Tabelle und die dazugehörigen Dokumente aus dem Google-Drive-Verzeichnis zusammen. Die Dokumente werden ein zweites mal auf ihre Korrektheit und Vollständigkeit überprüft und ebenso die übernommenen Daten in der Excel Tabelle. Die/der Zweitprüfende erstellt daraufhin mithilfe des Tools Ultradox eine Bewilligung oder Ablehnung. 

Handelt es sich um eine Ablehnung, wird diese der/dem Studierenden direkt per E-Mail zugesandt. Handelt es sich um eine Bewilligung, wird diese an das Finanzteam des \gls{asta} weitergeleitet um den Auszahlungsprozess anzustoßen. Sobald eine Rückmeldung des Finanzteams kommt, dass alles in Ordnung ist, wird die Bewilligung an die/den Studierenden versendet. 
- Prozessablauf+Visualisierung durch BPMN

\section{Probleme und Schwachstellen}\label{sec:Probleme_und_Schwachstellen}

\section{Übergang zum SOLL-Zustand}\label{sec:Übergang_Sollzustand}
+++ optional +++